% related_work.tex — Section 2 (Related Work), included from main.tex via % related_work.tex — Section 2 (Related Work), included from main.tex via % related_work.tex — Section 2 (Related Work), included from main.tex via % related_work.tex — Section 2 (Related Work), included from main.tex via \input{related_work}.
% Source: knowledge/50-publications/paper-1-biblio.md, section "Related Work (ébauche en anglais)".
% Citations use natbib (\citet / \citep); keys live in references.bib.

\paragraph{Backtest overfitting and the false-positive toolkit.}
A large literature protects strategy research against selection luck. \citet{bailey2012,bailey2014a} derive the probabilistic Sharpe ratio and the expected maximum of $N$ Sharpe ratios under the null, and deflate reported Sharpe ratios accordingly; \citet{bailey2014b,bailey2016} bound the minimum backtest length and estimate the probability of backtest overfitting through combinatorially symmetric cross-validation, later extended to purged and combinatorial cross-validation \citep{lopezdeprado2018}. In the multiple-testing tradition, \citet{white2000} and \citet{hansen2005} test the best of many models against the universe it was drawn from, \citet{harvey2016} raise the factor hurdle to $t > 3$, \citet{harvey2015} derive non-linear Sharpe haircuts, and \citet{chordia2020} calibrate thresholds on two million random strategies. \citet{bajgrowicz2012} and \citet{novymarx2015} apply the same logic to technical rules and combined signals. All of these instruments control the probability of adopting a strategy that has no signal. None of them reports the probability of rejecting a strategy that has one.

\paragraph{Power and missed discoveries.}
The false-negative side has been studied on a different object. \citet{harvey2020a} calibrate Type I and Type II errors jointly for the cross-section of fund alphas and conclude that existing methods lack power to detect outperforming managers; \citet{andrikogiannopoulou2019} simulate funds with known alphas and show that the false-discovery-rate procedure of \citet{barras2010} misclassifies up to 65\% of economically large alphas as zero. \citet{harvey2020b} compare the size and power of multiple-testing corrections by simulation, and \citet{inoue2005} show that out-of-sample tests sacrifice power relative to in-sample tests of the same hypothesis. These studies treat power as a property of a test applied to a distribution of effects, not of an operating pipeline with a decision frequency, a risk budget and a finite horizon.

\paragraph{Planted signals as positive controls.}
Validating a pipeline by injecting a known effect is standard outside finance: spike-in standards of known concentration measure the sensitivity of sequencing pipelines \citep{jiang2011}; synthetic positive controls, obtained by injecting simulated outcomes into real observational data, calibrate confidence intervals of epidemiological estimators \citep{schuemie2018}; active-control trials are uninterpretable without assay sensitivity \citep{temple2000}; randomization tests check that an explanation method reacts to the model it claims to explain \citep{adebayo2018}. Closest to our design, \citet{adeika2026} inject class-dependent effects of known amplitude into real EEG recordings and report a measured \emph{detection floor} per decoder, recommending that every null result be accompanied by one. The idea reached quantitative finance only recently. \citet{zheng2026} plants a constant alpha among candidate strategies and measures the adoption rate of admission gates built on the deflated Sharpe ratio against a random-admission control: gates remove false discoveries for weak signals at the cost of 93--99\% abstention, and criteria computed on absolute returns silently reject every true signal. \citet{liu2026} and \citet{qu2026} use planted-signal worlds to measure the size and power of LLM research agents and the admission delay of an anytime-valid referee. \citet{santoni2026} and \citet{gencay2026} evaluate composite robustness scores and search-aware deflation, again on the false-positive side.

\paragraph{Effective sample size, sequential stopping and censoring.}
Overlapping observations shrink the effective sample far below the nominal count \citep{hansen1980,boudoukh2008,lo2002}, and minimum-sample-size rules for prediction models are stated per event, not per row \citep{riley2019}. Sequential tests have a power function and an expected stopping time \citep{wald1945,doerks2026}, repeated looks inflate size \citep{armitage1969,lan1983,johari2022}, and a fixed fraction at risk determines the probability of ruin \citep{kelly1956,peters2019}. Incomplete observations are the domain of survival analysis \citep{kaplan1958}, and conditioning on survivors manufactures performance \citep{brown1992}.

\paragraph{Our contribution}
connects these strands: we treat the validation pipeline itself as the object whose power is measured, plant a signal of known directional strength aligned with realized outcomes, and report the smallest detectable signal per operating cell (decision cadence $\times$ risk per trade) and as a function of effective sample size, declaring a cell uninterpretable when the time budget censors the null or the signal. Unlike \citet{zheng2026}, we vary the operating cell rather than the gate, and find that the same pipeline on the same data is blind in some cells and detects a 54\% signal in others, and that a higher risk budget can make a strong true signal invisible.

% ---------------------------------------------------------------------------
% Citation check (2026-11-09): every \citet/\citep key above was matched against
% references.bib. Unresolved citations: none.
% ---------------------------------------------------------------------------
.
% Source: knowledge/50-publications/paper-1-biblio.md, section "Related Work (ébauche en anglais)".
% Citations use natbib (\citet / \citep); keys live in references.bib.

\paragraph{Backtest overfitting and the false-positive toolkit.}
A large literature protects strategy research against selection luck. \citet{bailey2012,bailey2014a} derive the probabilistic Sharpe ratio and the expected maximum of $N$ Sharpe ratios under the null, and deflate reported Sharpe ratios accordingly; \citet{bailey2014b,bailey2016} bound the minimum backtest length and estimate the probability of backtest overfitting through combinatorially symmetric cross-validation, later extended to purged and combinatorial cross-validation \citep{lopezdeprado2018}. In the multiple-testing tradition, \citet{white2000} and \citet{hansen2005} test the best of many models against the universe it was drawn from, \citet{harvey2016} raise the factor hurdle to $t > 3$, \citet{harvey2015} derive non-linear Sharpe haircuts, and \citet{chordia2020} calibrate thresholds on two million random strategies. \citet{bajgrowicz2012} and \citet{novymarx2015} apply the same logic to technical rules and combined signals. All of these instruments control the probability of adopting a strategy that has no signal. None of them reports the probability of rejecting a strategy that has one.

\paragraph{Power and missed discoveries.}
The false-negative side has been studied on a different object. \citet{harvey2020a} calibrate Type I and Type II errors jointly for the cross-section of fund alphas and conclude that existing methods lack power to detect outperforming managers; \citet{andrikogiannopoulou2019} simulate funds with known alphas and show that the false-discovery-rate procedure of \citet{barras2010} misclassifies up to 65\% of economically large alphas as zero. \citet{harvey2020b} compare the size and power of multiple-testing corrections by simulation, and \citet{inoue2005} show that out-of-sample tests sacrifice power relative to in-sample tests of the same hypothesis. These studies treat power as a property of a test applied to a distribution of effects, not of an operating pipeline with a decision frequency, a risk budget and a finite horizon.

\paragraph{Planted signals as positive controls.}
Validating a pipeline by injecting a known effect is standard outside finance: spike-in standards of known concentration measure the sensitivity of sequencing pipelines \citep{jiang2011}; synthetic positive controls, obtained by injecting simulated outcomes into real observational data, calibrate confidence intervals of epidemiological estimators \citep{schuemie2018}; active-control trials are uninterpretable without assay sensitivity \citep{temple2000}; randomization tests check that an explanation method reacts to the model it claims to explain \citep{adebayo2018}. Closest to our design, \citet{adeika2026} inject class-dependent effects of known amplitude into real EEG recordings and report a measured \emph{detection floor} per decoder, recommending that every null result be accompanied by one. The idea reached quantitative finance only recently. \citet{zheng2026} plants a constant alpha among candidate strategies and measures the adoption rate of admission gates built on the deflated Sharpe ratio against a random-admission control: gates remove false discoveries for weak signals at the cost of 93--99\% abstention, and criteria computed on absolute returns silently reject every true signal. \citet{liu2026} and \citet{qu2026} use planted-signal worlds to measure the size and power of LLM research agents and the admission delay of an anytime-valid referee. \citet{santoni2026} and \citet{gencay2026} evaluate composite robustness scores and search-aware deflation, again on the false-positive side.

\paragraph{Effective sample size, sequential stopping and censoring.}
Overlapping observations shrink the effective sample far below the nominal count \citep{hansen1980,boudoukh2008,lo2002}, and minimum-sample-size rules for prediction models are stated per event, not per row \citep{riley2019}. Sequential tests have a power function and an expected stopping time \citep{wald1945,doerks2026}, repeated looks inflate size \citep{armitage1969,lan1983,johari2022}, and a fixed fraction at risk determines the probability of ruin \citep{kelly1956,peters2019}. Incomplete observations are the domain of survival analysis \citep{kaplan1958}, and conditioning on survivors manufactures performance \citep{brown1992}.

\paragraph{Our contribution}
connects these strands: we treat the validation pipeline itself as the object whose power is measured, plant a signal of known directional strength aligned with realized outcomes, and report the smallest detectable signal per operating cell (decision cadence $\times$ risk per trade) and as a function of effective sample size, declaring a cell uninterpretable when the time budget censors the null or the signal. Unlike \citet{zheng2026}, we vary the operating cell rather than the gate, and find that the same pipeline on the same data is blind in some cells and detects a 54\% signal in others, and that a higher risk budget can make a strong true signal invisible.

% ---------------------------------------------------------------------------
% Citation check (2026-11-09): every \citet/\citep key above was matched against
% references.bib. Unresolved citations: none.
% ---------------------------------------------------------------------------
.
% Source: knowledge/50-publications/paper-1-biblio.md, section "Related Work (ébauche en anglais)".
% Citations use natbib (\citet / \citep); keys live in references.bib.

\paragraph{Backtest overfitting and the false-positive toolkit.}
A large literature protects strategy research against selection luck. \citet{bailey2012,bailey2014a} derive the probabilistic Sharpe ratio and the expected maximum of $N$ Sharpe ratios under the null, and deflate reported Sharpe ratios accordingly; \citet{bailey2014b,bailey2016} bound the minimum backtest length and estimate the probability of backtest overfitting through combinatorially symmetric cross-validation, later extended to purged and combinatorial cross-validation \citep{lopezdeprado2018}. In the multiple-testing tradition, \citet{white2000} and \citet{hansen2005} test the best of many models against the universe it was drawn from, \citet{harvey2016} raise the factor hurdle to $t > 3$, \citet{harvey2015} derive non-linear Sharpe haircuts, and \citet{chordia2020} calibrate thresholds on two million random strategies. \citet{bajgrowicz2012} and \citet{novymarx2015} apply the same logic to technical rules and combined signals. All of these instruments control the probability of adopting a strategy that has no signal. None of them reports the probability of rejecting a strategy that has one.

\paragraph{Power and missed discoveries.}
The false-negative side has been studied on a different object. \citet{harvey2020a} calibrate Type I and Type II errors jointly for the cross-section of fund alphas and conclude that existing methods lack power to detect outperforming managers; \citet{andrikogiannopoulou2019} simulate funds with known alphas and show that the false-discovery-rate procedure of \citet{barras2010} misclassifies up to 65\% of economically large alphas as zero. \citet{harvey2020b} compare the size and power of multiple-testing corrections by simulation, and \citet{inoue2005} show that out-of-sample tests sacrifice power relative to in-sample tests of the same hypothesis. These studies treat power as a property of a test applied to a distribution of effects, not of an operating pipeline with a decision frequency, a risk budget and a finite horizon.

\paragraph{Planted signals as positive controls.}
Validating a pipeline by injecting a known effect is standard outside finance: spike-in standards of known concentration measure the sensitivity of sequencing pipelines \citep{jiang2011}; synthetic positive controls, obtained by injecting simulated outcomes into real observational data, calibrate confidence intervals of epidemiological estimators \citep{schuemie2018}; active-control trials are uninterpretable without assay sensitivity \citep{temple2000}; randomization tests check that an explanation method reacts to the model it claims to explain \citep{adebayo2018}. Closest to our design, \citet{adeika2026} inject class-dependent effects of known amplitude into real EEG recordings and report a measured \emph{detection floor} per decoder, recommending that every null result be accompanied by one. The idea reached quantitative finance only recently. \citet{zheng2026} plants a constant alpha among candidate strategies and measures the adoption rate of admission gates built on the deflated Sharpe ratio against a random-admission control: gates remove false discoveries for weak signals at the cost of 93--99\% abstention, and criteria computed on absolute returns silently reject every true signal. \citet{liu2026} and \citet{qu2026} use planted-signal worlds to measure the size and power of LLM research agents and the admission delay of an anytime-valid referee. \citet{santoni2026} and \citet{gencay2026} evaluate composite robustness scores and search-aware deflation, again on the false-positive side.

\paragraph{Effective sample size, sequential stopping and censoring.}
Overlapping observations shrink the effective sample far below the nominal count \citep{hansen1980,boudoukh2008,lo2002}, and minimum-sample-size rules for prediction models are stated per event, not per row \citep{riley2019}. Sequential tests have a power function and an expected stopping time \citep{wald1945,doerks2026}, repeated looks inflate size \citep{armitage1969,lan1983,johari2022}, and a fixed fraction at risk determines the probability of ruin \citep{kelly1956,peters2019}. Incomplete observations are the domain of survival analysis \citep{kaplan1958}, and conditioning on survivors manufactures performance \citep{brown1992}.

\paragraph{Our contribution}
connects these strands: we treat the validation pipeline itself as the object whose power is measured, plant a signal of known directional strength aligned with realized outcomes, and report the smallest detectable signal per operating cell (decision cadence $\times$ risk per trade) and as a function of effective sample size, declaring a cell uninterpretable when the time budget censors the null or the signal. Unlike \citet{zheng2026}, we vary the operating cell rather than the gate, and find that the same pipeline on the same data is blind in some cells and detects a 54\% signal in others, and that a higher risk budget can make a strong true signal invisible.

% ---------------------------------------------------------------------------
% Citation check (2026-11-09): every \citet/\citep key above was matched against
% references.bib. Unresolved citations: none.
% ---------------------------------------------------------------------------
.
% Source: knowledge/50-publications/paper-1-biblio.md, section "Related Work (ébauche en anglais)".
% Citations use natbib (\citet / \citep); keys live in references.bib.

\paragraph{Backtest overfitting and the false-positive toolkit.}
A large literature protects strategy research against selection luck. \citet{bailey2012,bailey2014a} derive the probabilistic Sharpe ratio and the expected maximum of $N$ Sharpe ratios under the null, and deflate reported Sharpe ratios accordingly; \citet{bailey2014b,bailey2016} bound the minimum backtest length and estimate the probability of backtest overfitting through combinatorially symmetric cross-validation, later extended to purged and combinatorial cross-validation \citep{lopezdeprado2018}. In the multiple-testing tradition, \citet{white2000} and \citet{hansen2005} test the best of many models against the universe it was drawn from, \citet{harvey2016} raise the factor hurdle to $t > 3$, \citet{harvey2015} derive non-linear Sharpe haircuts, and \citet{chordia2020} calibrate thresholds on two million random strategies. \citet{bajgrowicz2012} and \citet{novymarx2015} apply the same logic to technical rules and combined signals. All of these instruments control the probability of adopting a strategy that has no signal. None of them reports the probability of rejecting a strategy that has one.

\paragraph{Power and missed discoveries.}
The false-negative side has been studied on a different object. \citet{harvey2020a} calibrate Type I and Type II errors jointly for the cross-section of fund alphas and conclude that existing methods lack power to detect outperforming managers; \citet{andrikogiannopoulou2019} simulate funds with known alphas and show that the false-discovery-rate procedure of \citet{barras2010} misclassifies up to 65\% of economically large alphas as zero. \citet{harvey2020b} compare the size and power of multiple-testing corrections by simulation, and \citet{inoue2005} show that out-of-sample tests sacrifice power relative to in-sample tests of the same hypothesis. These studies treat power as a property of a test applied to a distribution of effects, not of an operating pipeline with a decision frequency, a risk budget and a finite horizon.

\paragraph{Planted signals as positive controls.}
Validating a pipeline by injecting a known effect is standard outside finance: spike-in standards of known concentration measure the sensitivity of sequencing pipelines \citep{jiang2011}; synthetic positive controls, obtained by injecting simulated outcomes into real observational data, calibrate confidence intervals of epidemiological estimators \citep{schuemie2018}; active-control trials are uninterpretable without assay sensitivity \citep{temple2000}; randomization tests check that an explanation method reacts to the model it claims to explain \citep{adebayo2018}. Closest to our design, \citet{adeika2026} inject class-dependent effects of known amplitude into real EEG recordings and report a measured \emph{detection floor} per decoder, recommending that every null result be accompanied by one. The idea reached quantitative finance only recently. \citet{zheng2026} plants a constant alpha among candidate strategies and measures the adoption rate of admission gates built on the deflated Sharpe ratio against a random-admission control: gates remove false discoveries for weak signals at the cost of 93--99\% abstention, and criteria computed on absolute returns silently reject every true signal. \citet{liu2026} and \citet{qu2026} use planted-signal worlds to measure the size and power of LLM research agents and the admission delay of an anytime-valid referee. \citet{santoni2026} and \citet{gencay2026} evaluate composite robustness scores and search-aware deflation, again on the false-positive side.

\paragraph{Effective sample size, sequential stopping and censoring.}
Overlapping observations shrink the effective sample far below the nominal count \citep{hansen1980,boudoukh2008,lo2002}, and minimum-sample-size rules for prediction models are stated per event, not per row \citep{riley2019}. Sequential tests have a power function and an expected stopping time \citep{wald1945,doerks2026}, repeated looks inflate size \citep{armitage1969,lan1983,johari2022}, and a fixed fraction at risk determines the probability of ruin \citep{kelly1956,peters2019}. Incomplete observations are the domain of survival analysis \citep{kaplan1958}, and conditioning on survivors manufactures performance \citep{brown1992}.

\paragraph{Our contribution}
connects these strands: we treat the validation pipeline itself as the object whose power is measured, plant a signal of known directional strength aligned with realized outcomes, and report the smallest detectable signal per operating cell (decision cadence $\times$ risk per trade) and as a function of effective sample size, declaring a cell uninterpretable when the time budget censors the null or the signal. Unlike \citet{zheng2026}, we vary the operating cell rather than the gate, and find that the same pipeline on the same data is blind in some cells and detects a 54\% signal in others, and that a higher risk budget can make a strong true signal invisible.

% ---------------------------------------------------------------------------
% Citation check (2026-11-09): every \citet/\citep key above was matched against
% references.bib. Unresolved citations: none.
% ---------------------------------------------------------------------------
